\section{Terraform}
\textbf{Declarative vs. Imperative:} Terraform is \textit{declarative} -- you describe the desired end state and TF figures out the steps to get there. \textit{Imperative} (shell scripts, manual CLI calls) instead describes \textit{how} to reach a result step by step. \\
TF doesn't speak directly with an SDK, but rather Terraform -> Provider -> Client SDK. Diffrent providers enable diffrent platforms (AWS, Azure, Kubernetes, ...). \textbf{Resources} (\texttt{resource}) are the "what" TF creates and manages (e.g. an EC2 instance). \textbf{Data Sources} (\texttt{data}) are read-only lookups of existing infrastructure used as input (e.g. an existing AMI or VPC). A sample in HCL (Hashicorp Configuration Language):
\begin{lstlisting}[language=hcl]
variable "instance_type" {
  default = "t2.micro"
}
provider "aws" {
  region = "us-east-1"
}
resource "aws_instance" "web" {
  ami           = "ami-0c55b159cbfafe1f0"
  instance_type = var.instance_type
}
output "public_ip" {
  value = aws_instance.web.public_ip
}
\end{lstlisting}
To deploy infrastructure, write HCL in files like \texttt{main.tf}, then run \texttt{terraform init}, \texttt{terraform plan} to show changes that would be made, \texttt{terraform apply} to actually apply the changes. Use \texttt{terraform destroy} to delete all made changes.
\subsection{State}
Terraform stores state in the \texttt{terraform.tfstate} file, mapping desired (HCL) to actual (real) infrastructure. When working in teams, this state file also has to be shared as the terraform command relies on is validity. \\
\textbf{Remote backend:} store state remotely (e.g. an S3 Bucket) so the whole team shares it. \textbf{State locking} prevents two people running \texttt{apply} at once and corrupting the state; the S3 backend locks natively via \texttt{use\_lockfile} (older setups used a separate DynamoDB table).
\begin{lstlisting}[language=hcl]
terraform {
  backend "s3" {
    bucket         = "my-tfstate"
    key            = "prod/terraform.tfstate"
    region         = "us-east-1"
    use_lockfile   = true # S3-native state locking (replaces DynamoDB)
  }
}
\end{lstlisting}

\subsection{Modules}
Modules are reusable, parameterized templates (a folder of \texttt{.tf} files) to standardize and DRY up infrastructure. The root config is itself a module and can call child modules:
\begin{lstlisting}[language=hcl]
module "vpc" {
  source = "./modules/vpc" # or registry/git URL
  cidr   = "10.0.0.0/16"   # input variable
}
# reference an output: module.vpc.vpc_id
\end{lstlisting}

\subsection{More Commands \& Patterns}
\texttt{init} does three things: configure the backend, download providers, download modules. \texttt{fmt} formats, \texttt{validate} checks, \texttt{console} opens a REPL (print with \texttt{var.x}). \texttt{state list} lists tracked resources, \texttt{state show <addr>} shows their attributes. Reference any attribute as \texttt{<type>.<name>.<attr>} (e.g. \texttt{aws\_instance.web.public\_ip}).
\begin{lstlisting}[language=hcl]
import {                # adopt existing infra into state
  to = aws_instance.web
  id = "i-0abc123"      # provider-side ID; then: apply
}                       # CLI: terraform import aws_instance.web i-0abc123

variable "AMIS" {       # AMIs differ per region -> use a map
  type    = map(string)
  default = { us-east-1 = "ami-05b1..." }
}
ami = var.AMIS[var.region]   # select AMI for the region
\end{lstlisting}
